\documentclass[11pt,a4paper]{article}
\usepackage[T1]{fontenc}
\usepackage[utf8]{inputenc}
\usepackage{mathptmx}
\AtBeginDocument{\let\hbar\relax
  \DeclareMathSymbol{\hbar}{\mathord}{AMSb}{"7E}}
\usepackage{amsmath,amssymb,amsthm}
\usepackage{geometry}
\usepackage{microtype}
\usepackage{hyperref}
\usepackage{booktabs}
\usepackage{enumitem}
\usepackage[numbers]{natbib}
\usepackage{titlesec}

\geometry{top=1.1in,bottom=1.1in,left=1.15in,right=1.15in}

\titleformat{\section}{\normalfont\bfseries\large}{\thesection}{1em}{}
\titleformat{\subsection}{\normalfont\bfseries\normalsize}{\thesubsection}{1em}{}
\titleformat{\subsubsection}{\normalfont\bfseries\itshape\normalsize}
  {\thesubsubsection}{1em}{}
\titlespacing*{\section}{0pt}{14pt}{4pt}
\titlespacing*{\subsection}{0pt}{10pt}{3pt}

\theoremstyle{plain}
\newtheorem{theorem}{Theorem}[section]
\newtheorem{proposition}[theorem]{Proposition}
\newtheorem{corollary}[theorem]{Corollary}
\theoremstyle{definition}
\newtheorem{definition}[theorem]{Definition}
\newtheorem{prediction}[theorem]{Prediction}
\newtheorem{postulate}[theorem]{Postulate}
\theoremstyle{remark}
\newtheorem{remark}[theorem]{Remark}

\hypersetup{colorlinks=true,linkcolor=black,citecolor=black,urlcolor=black,
  pdftitle={Relational Actualism: Irreversible Events as Primitive Basis},
  pdfauthor={Joshua F. Sandeman}}

\newcommand{\Pact}{\mathcal{P}_{\mathrm{act}}}
\newcommand{\Tact}{T^{\mu\nu}_{\mathrm{RA}}}
\newcommand{\ARA}{A_{\mathrm{RA}}}
\newcommand{\ket}[1]{|#1\rangle}
\newcommand{\statusLV}{\hfill\textsc{[lv]}}
\newcommand{\statusCV}{\hfill\textsc{[cv]}}
\newcommand{\statusDR}{\hfill\textsc{[dr]}}
\newcommand{\statusAR}{\hfill\textsc{[ar]}}

% ─────────────────────────────────────────────────────────────────────────────
\begin{document}

% ── Title block ───────────────────────────────────────────────────────────────
\begin{center}
  {\LARGE\bfseries Relational Actualism:\par}
  \vspace{4pt}
  {\large\bfseries Irreversible Events as Primitive Basis\\
   of Physical Reality\par}
  \vspace{6pt}
  {\normalsize Paper~V of the Relational Actualism Series\par}
  \vspace{14pt}
  {\large Joshua F.\ Sandeman$^{1*}$\par}
  \vspace{4pt}
  {\normalsize $^{1}$Independent Researcher, Salem, OR, USA.\par}
  {\normalsize $^{*}$Corresponding author: sansuikyo@gmail.com\par}
  \vspace{14pt}
\end{center}

% ── Abstract ─────────────────────────────────────────────────────────────────
\begin{abstract}
This paper presents the complete logical structure of Relational Actualism~(RA),
a theoretical framework founded on a single ontological commitment: irreversible
actualization events---physical interactions that produce a permanent increase in
quantum relative entropy above a derived threshold---are the primitive basis of
physical reality.  From this commitment, implemented as an irreversibly growing
directed acyclic graph with a unique self-consistent growth rule (the BDG action
with coefficients $(1,-1,9,-16,8)$), the framework derives: the electromagnetic
and strong coupling constants, the Standard Model particle spectrum and gauge
group structure, the Koide lepton mass formula, the structural zero cosmological
constant, the dark matter mechanism, the Hubble tension resolution, the CMB
anomalies as signatures of parent SMBH dynamics, the dissolution of the measurement
problem, substrate-independent conditions for biological complexity, and the fate
of the universe as fragmentation rather than heat death.

The theory's logical structure is not a deduction chain but a \emph{constraint
web}: results reinforce each other through multiple independent paths.  This
paper maps the complete web, catalogues all predictions with their experimental
status, presents the formal verification record (163~Lean~4 theorems compiled
clean in Lean~4.29, one mechanical \texttt{sorry} tag requiring an external
library dependency), and identifies the remaining open problems---all technical
rather than foundational.  We argue that RA constitutes a falsifiable,
parameter-free account of physical reality from the Planck scale to the origin
of life.
\end{abstract}

%================================================================
\section{Introduction}\label{sec:intro}
%================================================================

Modern physics rests on two foundational frameworks---quantum mechanics and
general relativity---that are individually spectacularly successful and jointly
incompatible.  Quantum mechanics treats spacetime as a fixed background; general
relativity treats it as a dynamical variable.  The measurement problem asks how
quantum possibilities become classical facts.  The cosmological constant problem
asks why the vacuum energy is $10^{120}$ times smaller than predicted.  The dark
matter problem asks what constitutes 27\% of the universe's energy density.
These are not independent puzzles awaiting independent solutions.  They are
symptoms of a single conceptual gap.

Relational Actualism~(RA) addresses this gap with a single ontological
commitment: that irreversible actualization events are the primitive physical
reality.  Spacetime, conservation laws, quantum mechanics, and general
relativity are all emergent descriptions of the same underlying process---the
irreversible growth of a directed acyclic graph.

The companion papers in this series develop specific consequences: the unique
growth rule and its physical outputs (Paper~I~\cite{P1}), the dissolution of
the measurement problem (Paper~II~\cite{P2}), cosmological consequences from
nucleation to fragmentation (Paper~III~\cite{P3}), and substrate-independent
conditions for biological complexity (Paper~IV~\cite{P4}).  This paper presents
the framework as a whole: the logical architecture, the complete prediction
catalogue, the formal verification record, and the remaining open problems.

\subsection{Epistemological framing}\label{sec:epistemology}

The epistemological structure is stated plainly.  The growing DAG is
\emph{posited}.  The self-consistency constraints---locality, causal invariance,
non-redundancy, self-sustaining selectivity---are \emph{derived}.  The physical
quantities that follow (coupling constants, particle spectrum, cosmological
parameters, complexity thresholds) are \emph{predictions} of the posit, tested
against experiment.  Each agreement is evidence for the hypothesis.  A single
disagreement would falsify it.

RA does not ask for acceptance.  It asks for serious investigation.  The
framework makes specific, quantitative, falsifiable predictions across an
unusually wide range of experimental domains.  The purpose of this paper is to
make those predictions precise and to map the logical structure that connects
them.

%================================================================
\section{The Ontological Commitment}\label{sec:ontology}
%================================================================

RA's ontology consists of exactly one kind of entity: actualization events.
These are irreversible physical interactions that write permanent vertices and
directed edges into a growing directed acyclic graph~(DAG).  The vertices are
events.  The edges point from cause to effect.  Acyclicity---the impossibility
of causal loops---encodes the irreversibility of time at the most fundamental
level.

\paragraph{Time.}
Time is not a background dimension.  Time is the growing of the graph.  The past
is what has been written.  The future is what has not.  The present is the
process of writing the next vertex.

\paragraph{Space.}
Space is the network of actualized causal relations.  The metric---the notion of
distance---emerges from the density and structure of the graph at macroscopic
scales.  Proper time along a worldline is the integer count of actualization
events.

\paragraph{Conservation.}
Conservation is the Local Ledger Condition~(LLC): at every vertex, what flows in
equals what flows out.  This is not a law imposed on the graph.  It is the
graph's self-consistency condition, and it has been formally verified in
Lean~4~[LV].

\paragraph{The growth rule.}
The growth rule determines which proposed new vertices are accepted.  Paper~I
establishes that self-consistency admits exactly one growth rule: the BDG action
with coefficients $(1,-1,9,-16,8)$ in $d=4$ spacetime dimensions.  No other
coefficients and no other dimension are compatible with locality, causal
invariance, non-redundant information extraction, and self-sustaining
selectivity.

%================================================================
\section{The Constraint Web}\label{sec:web}
%================================================================

The logical structure of RA is not a deduction chain---a sequence of premises
leading to conclusions.  It is a \emph{constraint web}: a network of mutually
reinforcing results, each derivable from multiple independent paths through the
web.  This is the theory's most distinctive structural feature and its strongest
evidence for correctness.

\subsection{Why a web, not a chain}\label{sec:web-chain}

In a deduction chain, if one link breaks, everything downstream is lost.  In a
constraint web, the failure of any single connection leaves the remaining
connections intact.  The web's self-consistency provides a qualitatively
different kind of evidence from a linear derivation: the probability that
\emph{all} connections are simultaneously wrong by coincidence decreases
combinatorially with the number of independent paths.

RA's web has approximately nine interlocking constraint loops.

\subsection{The nine loops}\label{sec:loops}

\textbf{Loop~1 (BDG $\to$ couplings $\to$ spectrum).}
The BDG integers determine both the coupling constants and the particle
spectrum.  These are independently measured quantities that must be mutually
consistent.

\textbf{Loop~2 (LLC $\to$ GR).}
The Local Ledger Condition gives local conservation at every vertex (L01,~LV).
Amplitude locality~(O01,~LV) and $d=4$~(L11,~LV) together with the
Benincasa--Dowker uniqueness theorem~\cite{BenincasaDowker2010} determine the
BDG action as the unique local causal-set action, whose continuum limit is the
Einstein--Hilbert action.  The Bianchi identity is not a separate geometric
fact---it is the LLC viewed at the level of the metric.  General relativity is
not assumed; it is what the LLC \emph{looks like} in the continuum approximation.

\textbf{Loop~3 ($\Delta S^* \to$ measurement $\to$ Born rule).}
The actualization threshold $\Delta S^* \approx 0.601$~nats governs quantum
measurement and is consistent with the Born rule through the quantum measure
framework.

\textbf{Loop~4 ($\mu = 1 \to$ five scales).}
The Erd\H{o}s--R\'enyi percolation threshold at $\mu = 1$ appears in QCD
confinement, galactic rotation curves, quantum computing fault tolerance, the
actualization selectivity ceiling, and the origin of biological complexity.

\textbf{Loop~5 ($\Pact \to \Lambda = 0$, WIMP prohibition).}
The actualization projector~$\Pact$ removes vacuum contributions from the metric
source (giving $\Lambda = 0$ structurally) and excludes non-interacting
particles from gravitational coupling (giving the WIMP prohibition).

\textbf{Loop~6 (severance $\to$ nucleation $\to$ CMB).}
Causal severance at black hole horizons produces daughter universe nucleation.
The parent's Kerr geometry imprints a coherent quadrupole and octupole on the
daughter's initial conditions, explaining five CMB anomalies with three
parameters.

\textbf{Loop~7 (O14 $\to$ uniqueness).}
The Yeats moments plus M\"obius inversion prove the BDG coefficients are the
\emph{unique} solution to the self-consistency constraints.  This closes the
framework: there is no landscape of solutions.

\textbf{Loop~8 (void--filament $\to$ fragmentation).}
Density-dependent expansion feedback drives filamentary cosmic web topology as a
dynamical attractor.  Voids undergo starvation severance ($\mu < 1$); dense
boundaries undergo daughter nucleation ($\mu > \mu_2$).  The universe fragments
rather than equilibrating, prohibiting heat death and Boltzmann Brain nucleation.

\textbf{Loop~9 (F1~+~F2 $\to$ Causal Firewall).}
The actualization threshold $\Delta S^*$ and the $\mu = 1$ percolation transition
give substrate-independent conditions for biological complexity: redundant
information recording~(F1) and controlled actualization rate~(F2).

\subsection{Mutual implication of laws and phenomena}\label{sec:mutual}

The constraint web reveals that the traditional distinction between ``laws'' and
``phenomena'' is an artifact of how physics is presented, not a feature of the
underlying structure.  Conservation (the LLC) is not a law that actualization
events obey---it is what actualization events \emph{are}, viewed at the level of
bookkeeping.  The arrow of time is not explained by the framework---it is
encoded in the framework's most primitive commitment.  The coupling constants
are not inputs to the growth rule---they are outputs of it.

Every ``law'' in RA is a pattern that the web produces; every ``phenomenon'' is
an instance of a pattern.  The web dissolves the hierarchy by making both
derivable from the same source: the self-consistency of the growing DAG.

%================================================================
\section{Formal Verification Record}\label{sec:lean}
%================================================================

The discrete mathematical core of the framework has been formally verified in
Lean~4.29/Mathlib---the same proof assistant used to verify portions of Fermat's
Last Theorem.  Seven files compile with zero errors and zero warnings on core
results.

\begin{table}[ht]
\centering
\small
\begin{tabular}{@{}lrrl@{}}
\toprule
File & Theorems & Sorry & Key results \\
\midrule
\texttt{RA\_D1\_Proofs}       & 76 & 0 & L08 ($\alpha_{\mathrm{EM}}^{-1}=137$), L09 (Koide), L10--L12 (SM spectrum) \\
\texttt{RA\_GraphCore}         &  5 & 0 & L01 (LLC), L02 (Graph Cut), L03 (Markov Blanket) \\
\texttt{RA\_Koide}             & 16 & 0 & L09 ($K = 2/3$ for all $\theta$, all $m_0 > 0$) \\
\texttt{RA\_AmpLocality}       &  5 & 0 & O01 (amplitude locality), O02 (causal invariance) \\
\texttt{RA\_AQFT\_Proofs\_v10} &  9 & 1 & L04--L07 (frame independence, Rindler, CPTP) \\
\texttt{RA\_O14\_Uniqueness}   & 50 & 0 & O14 (BDG uniqueness via Yeats--Rota) \\
\texttt{RA\_Spin2\_Macro}      &  2 & 0 & O10s (massless spin-2 DOF count) \\
\midrule
\textbf{Total} & \textbf{163} & \textbf{1} & \\
\bottomrule
\end{tabular}
\caption{Lean~4 formal verification inventory (April~2026).  The single
\texttt{sorry} tag is a mechanical adapter requiring the Lean-QuantumInfo
external library; it is not a conceptual gap.  All files compile clean in
Lean~4.29.0 with current Mathlib4.}
\label{tab:lean}
\end{table}

The scope of these results is correctly characterised: they validate the
internal mathematical consistency of the formal claims---discrete graph
topology, finite-dimensional matrix algebra, relative entropy covariance,
integer arithmetic identities.  They do not prove that the universe executes the
BDG growth rule.  The physical claim rests on the predictions, not the Lean
proofs.

The amplitude locality theorem~(O01) deserves special emphasis.  It establishes
that the BDG amplitude function depends only on the causal past---a result that
follows from the transitivity of the causal order and the structure of the BDG
action, with no appeal to QFT microcausality or the continuum limit.  Combined
with the commutativity of complex multiplication, this gives unconditional
causal invariance~(O02): the quantum measure is independent of the choice of
linear extension.

%================================================================
\section{Derived Quantities}\label{sec:quantities}
%================================================================

The following quantities are derived from the BDG integers $(1,-1,9,-16,8)$ with
no free parameters.  Each is independently testable against experiment.

\begin{description}[style=nextline,leftmargin=1.5em]

\item[N01: Fine structure constant \statusLV]
$\alpha_{\mathrm{EM}}^{-1} = (137 + \sqrt{137^2 + 4 P_{\mathrm{acc}}\, c_2}\,)/2
= 137.036019$.  PDG value: $137.035999$.  Agreement: $0.00001\%$.
Every input---$137$~(LV), $P_{\mathrm{acc}}$~(CV), $c_2$~(LV)---is computable
from the BDG integers.

\item[N02: Strong coupling constant \statusCV]
$\alpha_s(m_Z) = 1/\sqrt{c_2\, c_4} = 1/\sqrt{72} \approx 0.11785$.
PDG value: $0.1180 \pm 0.0009$.  Agreement: $0.13\%$.

\item[N03: Actualization threshold \statusCV]
$\Delta S^* = -\ln P_{\mathrm{acc}}(1) = 0.60069$~nats (exact irrational).

\item[N04: Spin-bath collapse timescale \statusCV]
$t^* = (1/g)\sqrt{\Delta S^*/2} \approx 0.274/g$ (parameter-free).

\item[N05: Actualization density parameter \statusCV]
$\xi \approx 0.601\,\ell_P^{-2}$ (provisional).

\item[N06: Eridanus Hubble value \statusCV]
$H(\text{Eridanus}) \approx 75.9$~km/s/Mpc.

\item[N07: Proton mass \statusCV]
$m_p = \sqrt{c_4}\cdot\Lambda_{\mathrm{QCD}}$ to $0.08\%$.

\item[N08: Baryon-to-dark ratio \statusCV]
$f_0 = W_{\mathrm{other}}/W_{\mathrm{baryon}} = 17.32\times\alpha_s(m_p) = 5.42$.
Planck value: $5.416 \pm 0.015$.  Agreement: $0.3\%$.

\end{description}

%================================================================
\section{Complete Prediction Catalogue}\label{sec:predictions}
%================================================================

\subsection{Near-term (testable within 5~years)}\label{sec:pred-near}

\begin{prediction}[BMV null result]\label{pred:bmv}
No gravity-mediated entanglement at any mass scale.  The metric is sourced only
by actualized events; a mass in quantum superposition has $\ARA = 0$ and sources
zero gravitational field.  Multiple groups are pursuing this experiment.  A
positive result directly falsifies RA.
\end{prediction}

\begin{prediction}[Hubble gradient]\label{pred:hubble}
$H_0$ correlates linearly with baryon density $\rho_b$ along the line of sight.
Testable with DESI and large-scale structure surveys.
\end{prediction}

\begin{prediction}[WIMP floor]\label{pred:wimp}
No dark matter signal below the neutrino floor.  WIMPs, axions, and sterile
neutrinos are categorically excluded: any particle that does not participate in
on-shell actualization events has zero assembly depth and zero gravitational
coupling.
\end{prediction}

\begin{prediction}[Spin-bath timescale]\label{pred:spinbath}
$t^* = 0.274/g$, parameter-free.  Distinguishes RA from GRW/CSL collapse models
with different parameter dependences.  Testable with superconducting circuits.
\end{prediction}

\subsection{Medium-term (5--20~years)}\label{sec:pred-medium}

\begin{prediction}[Kinematic Coherence Bound]\label{pred:kcb}
$N_{\max} = \eta \cdot p_{\mathrm{th}}$ for fault-tolerant quantum arrays.  RA
predicts a hard ceiling; standard decoherence theory predicts exponentially
difficult but unbounded scaling.
\end{prediction}

\begin{prediction}[CMB axis persistence]\label{pred:cmb}
The quadrupole--octupole alignment persists in CMB-S4 at $> 3\sigma$.  The axis
direction correlates with large-scale cosmic web orientation (testable with
DESI/Euclid).  Five CMB anomalies follow from Kerr nucleation with three
parameters (Paper~III).
\end{prediction}

\begin{prediction}[Neutrino mass sum]\label{pred:neutrino}
$\sum m_\nu \approx 59$~meV (SU(3)$_{\mathrm{gen}}$ Koide fit).
CMB-S4 and Euclid will reach sensitivity $\sim$20~meV.
\end{prediction}

\subsection{Long-term (decades)}\label{sec:pred-long}

\begin{prediction}[Causal Firewall biosignatures]\label{pred:firewall}
Life is found wherever conditions F1 (redundant information recording) and F2
(controlled actualization rate) are satisfied, regardless of chemistry.  No
biosignatures where neither F1 nor F2 holds.
\end{prediction}

\begin{prediction}[No heat death]\label{pred:heatdeath}
The universe fragments via dual severance (void starvation + boundary
nucleation) before thermal equilibrium.  No single connected patch persists long
enough for Boltzmann Brain nucleation.
\end{prediction}

%================================================================
\section{Open Problems}\label{sec:open}
%================================================================

\subsection{Closed or superseded}\label{sec:open-closed}

\textbf{O02} (causal invariance without amplitude locality axiom): closed.
O01~(amplitude locality) is Lean-verified as a theorem of BDG
dynamics~\cite{Sandeman2026AmpLoc}.  O02 follows by
\texttt{List.Perm}~induction.

\textbf{O09} (covariant Step~4): closed.  Follows directly from O01 for the
unique BDG action.

\textbf{O10} (discrete Bianchi) and \textbf{O11} (Lorentz emergence):
dissolved.  The Bianchi identity is the LLC in continuum
costume~(\S\ref{sec:loops}, Loop~2).  Lorentz invariance is causal invariance in
continuum costume.  Neither requires a separate proof on the discrete structure;
both are what the DAG \emph{produces} when approximated as a smooth manifold.
The BDG action's continuum limit satisfies the contracted Bianchi
identity~\cite{BenincasaDowker2010} and exhibits Lorentz
invariance~\cite{DowkerGlaser2013}.

\subsection{Reframed and more tractable}\label{sec:open-reframed}

\textbf{O03} (Bianchi on curved backgrounds): reframed as convergence of the
known BDG operator on curved backgrounds.  Partial results
exist~\cite{MachetWang2020}.

\textbf{O06} ($\xi$ from first principles): essentially resolved.
$\xi = \Delta S^*/\ell_P^2$ with the unique discreteness scale.

\textbf{O12} (scale bridge): computational problem in the unique BDG dynamics,
no longer requiring a conceptual choice of dynamics.

\subsection{Genuinely open (technical)}\label{sec:open-tech}

\textbf{O04} (WEP formal bounds): equilibration dynamics calculation.

\textbf{O05} (Unruh detector model): open quantum systems calculation.

\textbf{O07} (BMV timescale): \emph{highest priority}.  Quantitative
decoherence criterion for macroscopic centre-of-mass position, needed before BMV
experimental data arrives.

\textbf{O08} (Causal Firewall $\tau_d$): first-principles derivation.
Partially addressed by non-Markovian corrections.

\textbf{O13} (fragmentation timescale): requires O12 first.

\medskip
\noindent\textbf{Net status:} zero foundational open problems.  Five technical
problems, of which one~(O07) is experimentally urgent.

%================================================================
\section{Relation to Other Programmes}\label{sec:other}
%================================================================

\subsection{Causal set theory}\label{sec:cst}

RA is built on the mathematical foundations of causal set
theory~\cite{BLMS1987,Sorkin2005}: the BDG action~\cite{BenincasaDowker2010},
the Rideout--Sorkin sequential growth dynamics~\cite{RideoutSorkin2000}, and the
hypothesis that spacetime is fundamentally discrete.  RA extends this programme
by~(i)~deriving the BDG coefficients intrinsically via the Yeats--Rota
uniqueness argument (Paper~I); (ii)~extracting physical observables (coupling
constants, particle spectrum, gauge groups) from the BDG integers; and
(iii)~taking the causal set as the fundamental ontology rather than an
approximation to a continuum manifold.

\subsection{Loop quantum gravity}\label{sec:lqg}

LQG discretises spacetime through spin networks and loop states.  RA discretises
through causal sets.  The key difference is that RA's discreteness is a property
of the \emph{growth process} (the DAG is built one vertex at a time), while
LQG's discreteness is a property of the \emph{state space} (area and volume
operators have discrete spectra).  RA makes specific predictions (coupling
constants, particle spectrum) that LQG does not.

\subsection{String theory}\label{sec:strings}

String theory introduces extra dimensions, supersymmetry, and a landscape of
vacua.  RA has no extra dimensions ($d = 4$ is uniquely selected), no
supersymmetry (the five topology types exhaust the SM without SUSY partners),
and no landscape (the growth rule is unique).  The two programmes make opposite
predictions about BSM physics: string theory predicts supersymmetric partners;
RA predicts none.

\subsection{Assembly Theory}\label{sec:at}

Assembly Theory~\cite{Sharma2023} measures molecular complexity via the assembly
index.  RA's assembly depth~$\tilde{A}_{\mathrm{RA}}$ is grounded in
fundamental physics (actualization events) rather than chemical operations.
Paper~IV shows the two converge for sufficiently complex molecules.

%================================================================
\section{The Democratisation of Foundational Physics}\label{sec:demo}
%================================================================

This framework was developed by an independent researcher working outside
academic institutions.  The tools that made this possible---formal proof
assistants (Lean~4), large language model collaboration (for derivation
execution, literature search, and red-team review), open-access preprint servers
(Zenodo), and computational algebra systems---are genuinely new.  They lower the
barriers to serious foundational physics research in a way that changes who can
contribute.

The framework is fully open.  All papers are available on Zenodo with permanent
DOIs.  The Lean~4 proof files are public.  The computation scripts are public.
Collaboration is actively sought---particularly from researchers in algebraic
QFT, causal set theory, analytic combinatorics, computational biology, and
quantum computing.

%================================================================
\section{Conclusions}\label{sec:conclusions}
%================================================================

Relational Actualism posits that irreversible actualization events, writing
permanent vertices into a growing directed acyclic graph, are the primitive
physical reality.  From this single commitment:

\begin{itemize}[nosep]
\item The growth rule is uniquely determined (Paper~I).
\item The coupling constants $\alpha_{\mathrm{EM}}$ and $\alpha_s$ are derived
  to $0.00001\%$ and $0.13\%$ (Paper~I).
\item The Standard Model particle spectrum is exhaustive and the gauge group
  $\mathrm{SU}(3)\times\mathrm{SU}(2)\times\mathrm{U}(1)$ follows from
  inclusion--exclusion counting (Paper~I).
\item The measurement problem is dissolved (Paper~II).
\item Three specific quantum predictions distinguish RA from all existing
  interpretations (Paper~II).
\item The cosmological constant is structurally zero (Paper~III).
\item The CMB axis anomaly and hemisphere asymmetry are signatures of parent
  SMBH dynamics at nucleation (Paper~III).
\item The universe fragments rather than equilibrates (Paper~III).
\item The origin of life is a phase transition at $\mu = 1$ (Paper~IV).
\item The five-scale $\mu = 1$ unification connects QCD to biology (Paper~IV).
\end{itemize}

The theory's logical structure is a constraint web, not a deduction chain.  Each
result is derivable from multiple independent paths, and the web's
self-consistency provides combinatorial evidence for correctness.  The formal
verification record (163~Lean~4 theorems) establishes the mathematical core
beyond reasonable doubt.

Zero foundational open problems remain.  Five technical problems await
resolution, of which one (O07, the BMV timescale) is experimentally urgent.  The
framework makes specific, falsifiable predictions across particle physics,
quantum mechanics, cosmology, quantum computing, and astrobiology.

Einstein asked whether God had any choice in making the universe.  The
uniqueness theorem answers: no.  Self-consistency admits exactly one solution.
The integers $(1,-1,9,-16,8)$ are not parameters.  The dimension $d = 4$ is not
a choice.  The coupling constants are not inputs.  They are the unique output of
a combinatorial identity---the only way for an irreversibly growing causal graph
to be self-consistent.

% ─────────────────────────────────────────────────────────────────────────────
\begin{thebibliography}{99}

\bibitem{P1}
J.~F.~Sandeman,
``Paper~I: The Physics of a Self-Consistent Causal Graph,''
Preprint (2026).

\bibitem{P2}
J.~F.~Sandeman,
``Paper~II: Wave Function Collapse as Irreversible Entropy Production,''
Preprint (2026).

\bibitem{P3}
J.~F.~Sandeman,
``Paper~III: The Origin, Structure, and Fate of Universes,''
Preprint (2026).

\bibitem{P4}
J.~F.~Sandeman,
``Paper~IV: The Causal Firewall,''
Preprint (2026).

\bibitem{P0}
J.~F.~Sandeman,
``Paper~0: A Biography of the Universe,''
Preprint (2026).

\bibitem{BenincasaDowker2010}
D.~M.~T.~Benincasa and F.~Dowker,
``The scalar curvature of a causal set,''
\textit{Phys.\ Rev.\ Lett.}\ \textbf{104}, 181301 (2010).

\bibitem{DowkerGlaser2013}
F.~Dowker and L.~Glaser,
``Causal set d'Alembertians for various dimensions,''
\textit{Class.\ Quantum Grav.}\ \textbf{30}, 195016 (2013).

\bibitem{MachetWang2020}
L.~Machet and J.~Wang,
``On the continuum limit of the causal set d'Alembertian in curved spacetimes,''
\textit{Class.\ Quantum Grav.}\ \textbf{38}, 015010 (2020).

\bibitem{BLMS1987}
L.~Bombelli, J.~Lee, D.~Meyer, and R.~D.~Sorkin,
``Space-time as a causal set,''
\textit{Phys.\ Rev.\ Lett.}\ \textbf{59}, 521 (1987).

\bibitem{Sorkin2005}
R.~D.~Sorkin,
``Causal sets: Discrete gravity,''
in \textit{Lectures on Quantum Gravity}
(Springer, 2005), pp.~305--327.

\bibitem{RideoutSorkin2000}
D.~Rideout and R.~D.~Sorkin,
``A classical sequential growth dynamics for causal sets,''
\textit{Phys.\ Rev.\ D} \textbf{61}, 024002 (2000).

\bibitem{Sharma2023}
M.~D.~Sharma \textit{et al.},
``Assembly theory explains and quantifies selection and evolution,''
\textit{Nature} \textbf{622}, 278--283 (2023).

\bibitem{Yeats2025}
K.~Yeats,
``Combinatorial interpretation of the coefficients of the causal set
d'Alembertian,''
\textit{Class.\ Quantum Grav.}\ \textbf{42}, 145003 (2025);
arXiv:2412.14036.

\bibitem{Rota1964}
G.-C.~Rota,
``On the foundations of combinatorial theory~I: Theory of M\"obius functions,''
\textit{Z.\ Wahrscheinlichkeitstheorie} \textbf{2}, 340--368 (1964).

\bibitem{Sandeman2026AmpLoc}
J.~F.~Sandeman,
``RA\_AmpLocality.lean: Amplitude locality as a theorem of BDG discrete DAG
dynamics,''
GitHub (2026), \url{github.com/jsandeman/Relational-Actualism}.

\end{thebibliography}

\end{document}
